% ==============================================================================
% CONCLUSION SECTION
% Path forward
% ==============================================================================

% ==============================================================================
% SLIDE 10
% ==============================================================================
\begin{frame}[t]
    \frametitle{Path forward: More empirical work and fleshing out the model}
    \small

    \begin{columns}[T,totalwidth=\textwidth]
        \begin{column}{0.50\textwidth}
            \begin{block}{Finish the empirical bridge}
                \begin{itemize}
                    \item Digitize historical council boundaries and old city council rosters
                    \item Use text analysis and additional data sources to measure member deference activity explicitly
                \end{itemize}
            \end{block}
        \end{column}
          \begin{column}{0.50\textwidth}
            \begin{block}{Flesh out the model}
                \begin{itemize}
                    \item Build out the spatial equilibrium: residents, entrants, rents, wages, commuting, amenities, and assets.
                    \item Add housing supply and approval stage: separate physical construction costs from the political approval wedge.
                    \item Define welfare objects by group: incumbent owners, incumbent renters, entrants, and aggregate welfare.
                \end{itemize}
            \end{block}
        \end{column}
    \end{columns}

    \vspace{0.01em}
    \begin{center}
        \textbf{Goal:} separate the value of local information from the cost of local veto power and examine the potential for compensation or pricing to mitigate the political approval wedge while retaining the value of local information.
    \end{center}
\end{frame}
